% !TEX program = pdflatex
\documentclass[11pt,a4paper]{article}

% =========================================================================
% GUÍA DE TRABAJO DEL INTEGRANTE
% Proyecto: CasaScore (AlquilaPuno)
% Gestión de Proyectos TICS (EST315) -- Semestre 2026-II -- Team Sobrevivientes
% =========================================================================

% --- PAQUETES DE CONFIGURACIÓN ---
\usepackage[utf8]{inputenc}
\usepackage[T1]{fontenc}
\usepackage[spanish,es-tabla]{babel}
\usepackage{geometry}
\geometry{
    top=3.6cm,
    bottom=2.3cm,
    left=2.2cm,
    right=2.2cm,
    % PARTE MODIFICABLE: headheight debe ser mayor que la altura del logo más alto (2.56cm = 73pt).
    headheight=80pt,
    headsep=15pt,
    footskip=20pt
}

\usepackage{graphicx}
\usepackage{fancyhdr}
\usepackage[table]{xcolor}
\usepackage{titlesec}
\usepackage{booktabs}
\usepackage{tabularx}
\usepackage{array}
\usepackage{amsmath}
\usepackage{enumitem}
\usepackage{microtype}
\usepackage{amssymb}
% --- RECURSOS VISUALES: GRÁFICOS Y CÓDIGO DE COLORES ---
\usepackage{tikz}
\usetikzlibrary{positioning,calc}
\definecolor{GrisFila}{RGB}{246, 248, 251}
\definecolor{CuadCerca}{RGB}{150, 30, 30}
\definecolor{CuadSatisf}{RGB}{180, 130, 20}
\definecolor{CuadInfo}{RGB}{10, 45, 90}
\definecolor{CuadMonit}{RGB}{110, 115, 125}
\newcommand{\etiqueta}[2]{{\setlength{\fboxsep}{2pt}\colorbox{#1}{\textcolor{white}{\scriptsize\bfseries #2}}}}
\newcommand{\cerca}{\etiqueta{CuadCerca}{Gestionar de cerca}}
\newcommand{\satisf}{\etiqueta{CuadSatisf}{Mantener satisfecho}}
\newcommand{\informado}{\etiqueta{CuadInfo}{Mantener informado}}
\newcommand{\monit}{\etiqueta{CuadMonit}{Monitorear}}
\usepackage{needspace}
\usepackage{longtable}
\definecolor{Rama1}{RGB}{10, 45, 90}
\definecolor{Rama2}{RGB}{40, 95, 150}
\definecolor{Rama3}{RGB}{180, 130, 20}
\definecolor{Rama4}{RGB}{30, 110, 60}
\definecolor{Rama5}{RGB}{110, 60, 140}
\definecolor{RaciR}{RGB}{150, 30, 30}
\definecolor{RaciA}{RGB}{10, 45, 90}
\definecolor{RaciC}{RGB}{180, 130, 20}
\definecolor{RaciI}{RGB}{140, 145, 155}
\definecolor{NivAlto}{RGB}{30, 110, 60}
\definecolor{NivMedio}{RGB}{180, 130, 20}
\definecolor{NivBajo}{RGB}{150, 30, 30}
\usetikzlibrary{patterns}
\definecolor{SevBaja}{RGB}{30, 110, 60}
\definecolor{SevMedia}{RGB}{180, 130, 20}
\definecolor{SevAlta}{RGB}{150, 30, 30}
\usetikzlibrary{arrows.meta,shapes.geometric,fit,backgrounds}
\usepackage{listings}
\lstset{basicstyle=\ttfamily\scriptsize, frame=single, rulecolor=\color{AzulUNA!40}, backgroundcolor=\color{GrisFila}, columns=fullflexible, keepspaces=true, breaklines=true, xleftmargin=4pt, xrightmargin=4pt, literate={á}{{\'a}}1 {é}{{\'e}}1 {í}{{\'i}}1 {ó}{{\'o}}1 {ú}{{\'u}}1 {ñ}{{\~n}}1 {Á}{{\'A}}1}

\usepackage{hyperref}

% Logos: funciona a 2 o 3 niveles de profundidad dentro de docs/latex/
\graphicspath{{../assets/img/}{../../assets/img/}{../../../assets/img/}}

% --- PALETA INSTITUCIONAL FIEI / UNA PUNO ---
\definecolor{AzulUNA}{RGB}{10, 45, 90}
\definecolor{DoradoFIEI}{RGB}{180, 130, 20}
\definecolor{GrisOscuro}{RGB}{50, 50, 50}
\definecolor{AzulClaro}{RGB}{228, 236, 246}

% --- CONFIGURACIÓN DE LISTAS ---
\setlist{topsep=2pt, itemsep=2pt, parsep=1pt}

% --- COLUMNAS DE TABLA ---
\newcolumntype{L}{>{\raggedright\arraybackslash}X}
\newcolumntype{P}[1]{>{\raggedright\arraybackslash}p{#1}}
\newcolumntype{C}[1]{>{\centering\arraybackslash}p{#1}}
\newcolumntype{Q}[1]{>{\raggedleft\arraybackslash}p{#1}}
\newcolumntype{R}{>{\raggedleft\arraybackslash}X}
\renewcommand{\arraystretch}{1.25}

% --- FORMATO DE ENLACES ---
\hypersetup{
    colorlinks=true,
    linkcolor=AzulUNA,
    urlcolor=AzulUNA,
    citecolor=AzulUNA,
    pdftitle={Guía De Trabajo Del Integrante -- CasaScore (AlquilaPuno)},
    pdfauthor={Team Sobrevivientes}
}

% --- FORMATO DE TÍTULOS ---
\titleformat{\section}{\color{AzulUNA}\normalfont\large\bfseries}{\thesection.}{0.7em}{}[\color{AzulUNA}\titlerule]
\titleformat{\subsection}{\color{AzulUNA}\normalfont\normalsize\bfseries}{\thesubsection.}{0.5em}{}
\titlespacing*{\section}{0pt}{9pt}{4pt}
\titlespacing*{\subsection}{0pt}{6pt}{2pt}

% --- ENCABEZADO Y PIE DE PÁGINA (fancyhdr) ---
\pagestyle{fancy}
\fancyhf{}

% ===== TAMAÑO DE LOS LOGOS (PARTE MODIFICABLE) =====
% logo_fiei.png  : height=1.45cm (lado izquierdo)
% logo_equipo.png: height=2.56cm (lado derecho; 25 % más grande porque la imagen tiene más borde blanco)
% Para cambiar el tamaño, modifique solo el valor height=... de cada \includegraphics.
% Si el logo más alto supera 2.6cm, aumente también headheight y top en la configuración de página.
\fancyhead[L]{%
    \includegraphics[height=1.45cm,keepaspectratio]{logo_fiei.png}%
}
\fancyhead[C]{%
    \scriptsize \textbf{\textcolor{AzulUNA}{UNIVERSIDAD NACIONAL DEL ALTIPLANO DE PUNO}}\\[2pt]
    \tiny \textcolor{GrisOscuro}{FACULTAD DE INGENIERÍA ESTADÍSTICA E INFORMÁTICA}\\[2pt]
    \scriptsize \textbf{\textcolor{AzulUNA}{Gestión de Proyectos TICS (EST315) \textbar\ Team Sobrevivientes}}\\[2pt]
    \footnotesize \textbf{\textcolor{DoradoFIEI}{GUÍA DE TRABAJO DEL INTEGRANTE}}%
}
\fancyhead[R]{%
    \includegraphics[height=2.56cm,keepaspectratio]{logo_equipo.png}%
}

\renewcommand{\headrulewidth}{0pt}
\renewcommand{\headrule}{}

\fancyfoot[L]{\scriptsize \textcolor{GrisOscuro}{CasaScore (AlquilaPuno)}}
\fancyfoot[C]{\small \thepage}
\fancyfoot[R]{\scriptsize \textcolor{GrisOscuro}{Puno, Perú -- Semestre 2026-II}}
\renewcommand{\footrulewidth}{0.4pt}

\begin{document}
\shorthandoff{<>}
\lstset{basicstyle=\ttfamily\fontsize{6.8}{8.2}\selectfont, breaklines=false}
\begin{center}
    {\Large \textbf{\textcolor{AzulUNA}{GUÍA DE TRABAJO DEL INTEGRANTE}}}\\[3pt]
    {\normalsize \textbf{Viviana Mary Quisbert Quispe -- Gestora de Comunicación, Riesgos y Documentación}}\\[4pt]
    \footnotesize \textbf{Fecha:} jueves 05/11/2026 (semana 11) \quad \textbar\quad \textbf{Versión:} 1.0 \quad \textbar\quad \textbf{Encarga:} Clyde N. Paricahua Pari
\end{center}

\vspace{0.1cm}

\section{Propósito}
El equipo trabaja desde ahora en el repositorio \texttt{ProyectoTics}. Esta guía explica, desde cero, cómo preparar tu computadora, clonar el repositorio y subir tus documentos desde la terminal, cada uno con su rama, su commit y su Pull Request. El seguimiento de tus tareas se hace en el navegador, en el tablero \textit{CasaScore -- Tablero híbrido}.

\vspace{0.1cm}
\noindent\begin{tabularx}{\textwidth}{@{}P{3.6cm} L@{}}
\toprule
\textbf{Repositorio:} & \nolinkurl{https://github.com/xClaIDx/ProyectoTics} \\
\textbf{Tablero:} & \nolinkurl{https://github.com/users/xClaIDx/projects/1} \\
\textbf{Tu usuario:} & \texttt{Viviana-11} \\
\textbf{Revisor:} & Clyde (\texttt{xClaIDx}) revisa, aprueba y fusiona cada Pull Request \\
\textbf{Plazo:} & PR 1 a 5 hasta el viernes 06/11; PR 6 el viernes 13/11, fecha del plan de cierre \\
\bottomrule
\end{tabularx}

\section{Lo que Recibes}
\begin{enumerate}[leftmargin=*]
    \item \textbf{La carpeta \texttt{paquete\_viviana}} (comprimida en un \texttt{.zip}), que Clyde te envía por WhatsApp. Contiene esta guía y la carpeta \nolinkurl{paquete_viviana/docs/latex/} con tus 6 documentos (\texttt{.tex} y \texttt{.pdf}) en las mismas carpetas que tienen en el repositorio.
    \item \textbf{La invitación al repositorio} \texttt{ProyectoTics}, que llega a tu correo de GitHub.
    \item \textbf{Tus tarjetas en el tablero}, ya asignadas a tu usuario.
\end{enumerate}

\Needspace{20\baselineskip}
\section{Preparación desde Cero (una sola vez)}

\subsection{Instalar Git y la terminal de GitHub}
{\small
\rowcolors{2}{white}{GrisFila}
\noindent\begin{tabularx}{\textwidth}{@{}P{2.2cm} L L@{}}
\toprule
\rowcolor{AzulClaro}
& \textbf{Windows} & \textbf{Mac} \\
\midrule
Git & Descarga \textbf{Git for Windows} desde \nolinkurl{https://git-scm.com/download/win} e instálalo con las opciones por defecto. & Abre la app \textbf{Terminal} y escribe \texttt{xcode-select --install}; acepta la instalación. \\
GitHub CLI & Descarga el instalador desde \nolinkurl{https://cli.github.com} e instálalo. & Descarga el instalador \texttt{.pkg} desde \nolinkurl{https://cli.github.com} e instálalo. \\
Terminal que usarás & \textbf{Git Bash} (búscalo en el menú Inicio). Todos los comandos de esta guía funcionan en Git Bash. & La app \textbf{Terminal}. \\
\bottomrule
\end{tabularx}
}

\vspace{0.1cm}
\noindent Cierra y vuelve a abrir la terminal, y comprueba que ambos programas responden con su número de versión:

\begin{lstlisting}
git --version
gh --version
\end{lstlisting}

\subsection{Configurar tu nombre}
Así cada commit queda a tu nombre. Usa el correo de tu cuenta de GitHub.

\begin{lstlisting}
git config --global user.name "Viviana Mary Quisbert Quispe"
git config --global user.email "tu-correo-de-github@ejemplo.com"
git config --global init.defaultBranch main
\end{lstlisting}

\Needspace{14\baselineskip}
\subsection{Iniciar sesión en GitHub desde la terminal}
\begin{lstlisting}
gh auth login
\end{lstlisting}

\noindent Responde así a cada pregunta (usa las flechas y Enter): \textbf{GitHub.com} $\rightarrow$ \textbf{HTTPS} $\rightarrow$ \textbf{Yes} (autenticar Git con tus credenciales) $\rightarrow$ \textbf{Login with a web browser}. La terminal muestra un código de 8 caracteres; pulsa Enter, se abre el navegador, pega el código y autoriza. Comprueba con:

\begin{lstlisting}
gh auth status
\end{lstlisting}

\Needspace{12\baselineskip}
\subsection{Aceptar la invitación al repositorio}
Desde el navegador, en el correo de GitHub (\textbf{View invitation} $\rightarrow$ \textbf{Accept invitation}) o en \nolinkurl{https://github.com/xClaIDx/ProyectoTics/invitations}. Sin este paso, la terminal rechazará tus envíos con un error de permiso.

\Needspace{16\baselineskip}
\subsection{Clonar el repositorio y preparar tu paquete}
\begin{enumerate}[leftmargin=*]
    \item Descomprime \texttt{paquete\_viviana.zip} con doble clic y deja la carpeta \texttt{paquete\_viviana} en tu \textbf{Escritorio}.
    \item En la terminal, ve al Escritorio y clona el repositorio:
\end{enumerate}

\begin{lstlisting}
cd ~/Desktop
git clone https://github.com/xClaIDx/ProyectoTics.git
cd ProyectoTics
ls docs/latex
\end{lstlisting}

\noindent Debes ver las carpetas \texttt{00\_comun}, \texttt{casascore} y \texttt{horus\_sentinel}. En Windows, si tu Escritorio está en OneDrive, usa \nolinkurl{cd ~/OneDrive/Escritorio} en lugar de \nolinkurl{cd ~/Desktop}. Desde ahora, todos los comandos se escriben \textbf{dentro de la carpeta} \texttt{ProyectoTics}; el paquete queda al lado, en \nolinkurl{../paquete_viviana}.


\Needspace{12\baselineskip}
\section{Cómo se Simula el Historial}
Cada documento conserva la fecha que tiene escrita; no se modifica ningún archivo. GitHub registrará el día real en que lo subes. Para que se sepa cuándo debió entregarse, el mensaje del commit y el título del Pull Request terminan con la fecha prevista, por ejemplo:

\begin{lstlisting}
docs(comun-conformacion): acta de conformacion del equipo [fecha prevista 01/09/2026]
\end{lstlisting}

\Needspace{20\baselineskip}
\section{Tus Seis Pull Requests}
Súbelos en este orden. Cada carpeta se copia desde \nolinkurl{../paquete_viviana/docs/latex/} a la misma ruta dentro de \nolinkurl{docs/latex/}.

\vspace{0.1cm}
{\small
\rowcolors{2}{white}{GrisFila}
\noindent\begin{tabularx}{\textwidth}{@{}C{0.6cm} P{4.6cm} L C{1.8cm}@{}}
\toprule
\rowcolor{AzulClaro}
\textbf{PR} & \textbf{Rama} & \textbf{Carpeta que subes} & \textbf{Fecha prevista} \\
\midrule
1 & \nolinkurl{docs/comun-conformacion} & \nolinkurl{00_comun/C2_acta_conformacion} & 01/09/2026 \\
2 & \nolinkurl{docs/comun-politicas} & \nolinkurl{00_comun/C1_politicas} & 04/09/2026 \\
3 & \nolinkurl{docs/casascore-riesgos} & \nolinkurl{casascore/10_plan_riesgos} & 22/09/2026 \\
4 & \nolinkurl{docs/casascore-comunicaciones} & \nolinkurl{casascore/12_plan_comunicaciones} & 22/09/2026 \\
5 & \nolinkurl{docs/casascore-actas-sprint} & \nolinkurl{casascore/15_actas_sprint} & 16/10/2026 \\
6 & \nolinkurl{docs/casascore-cierre} & \nolinkurl{casascore/23_plan_cierre_lecciones} & 13/11/2026 \\
\bottomrule
\end{tabularx}
}

\vspace{0.1cm}
\noindent El PR 6 corresponde a un documento con fecha 13/11/2026: súbelo ese día, no antes.

\Needspace{30\baselineskip}
\section{Ciclo Completo de un Pull Request}
Ejemplo con el PR 1. Cada bloque va precedido de lo que hace.

\vspace{0.1cm}
\noindent\textbf{1. Partir de la versión más reciente de \texttt{main}}
\begin{lstlisting}
git switch main
git pull
\end{lstlisting}

\noindent\textbf{2. Crear tu rama}
\begin{lstlisting}
git switch -c docs/comun-conformacion
\end{lstlisting}

\noindent\textbf{3. Copiar tu carpeta desde el paquete}
\begin{lstlisting}
cp -r ../paquete_viviana/docs/latex/00_comun/C2_acta_conformacion docs/latex/00_comun/
\end{lstlisting}

\noindent\textbf{4. Revisar qué cambió}: debe aparecer en rojo la carpeta nueva.
\begin{lstlisting}
git status
\end{lstlisting}

\noindent\textbf{5. Preparar y guardar el cambio (commit)}
\begin{lstlisting}
git add docs/latex/00_comun/C2_acta_conformacion
git commit -m "docs(comun-conformacion): acta de conformacion del equipo [fecha prevista 01/09/2026]"
\end{lstlisting}

\noindent\textbf{6. Enviar la rama a GitHub}
\begin{lstlisting}
git push -u origin docs/comun-conformacion
\end{lstlisting}

\noindent\textbf{7. Abrir el Pull Request con Clyde como revisor}
\begin{lstlisting}
gh pr create --base main --reviewer xClaIDx \
  --title "docs(comun-conformacion): acta de conformacion del equipo [fecha prevista 01/09/2026]" \
  --body "Fecha prevista 01/09/2026. Acta de conformación del equipo (.tex y .pdf). PR 1 de 6."
\end{lstlisting}

\noindent La terminal responde con el enlace del Pull Request. No lo fusiones tú: Clyde lo revisa y lo fusiona.

\vspace{0.1cm}
\noindent\textbf{8. Cuando Clyde lo fusione, volver a \texttt{main} y borrar la rama local}
\begin{lstlisting}
git switch main
git pull
git branch -d docs/comun-conformacion
\end{lstlisting}

\noindent No esperes a que fusione un PR para empezar el siguiente: el paso 1 siempre parte de \texttt{main}, y cada rama es independiente.

\Needspace{12\baselineskip}
\subsection{Si Clyde pide cambios}
Vuelve a tu rama, corrige el archivo, y envía un nuevo commit; el Pull Request se actualiza solo.
\begin{lstlisting}
git switch docs/comun-conformacion
git add docs/latex/00_comun/C2_acta_conformacion
git commit -m "fix(comun-conformacion): correcciones de la revision"
git push
\end{lstlisting}

\Needspace{12\baselineskip}
\subsection{Ver el estado de tus Pull Requests}
\begin{lstlisting}
gh pr status
gh pr list --author "@me"
gh pr view --web
\end{lstlisting}

\Needspace{16\baselineskip}
\section{Comandos de los PR 2 a 6}
Cada bloque se ejecuta igual que el ciclo anterior. Después de cada uno, cuando Clyde lo fusione, haz el paso 8 con el nombre de esa rama.

\Needspace{11\baselineskip}
\noindent\textbf{PR 2 -- Políticas del proyecto y gobernanza (04/09/2026)}
\begin{lstlisting}
git switch main && git pull
git switch -c docs/comun-politicas
cp -r ../paquete_viviana/docs/latex/00_comun/C1_politicas docs/latex/00_comun/
git add docs/latex/00_comun/C1_politicas
git commit -m "docs(comun-politicas): politicas del proyecto y gobernanza [fecha prevista 04/09/2026]"
git push -u origin docs/comun-politicas
gh pr create --base main --reviewer xClaIDx --fill
\end{lstlisting}

\Needspace{11\baselineskip}
\noindent\textbf{PR 3 -- Plan y registro de riesgos (22/09/2026)}
\begin{lstlisting}
git switch main && git pull
git switch -c docs/casascore-riesgos
cp -r ../paquete_viviana/docs/latex/casascore/10_plan_riesgos docs/latex/casascore/
git add docs/latex/casascore/10_plan_riesgos
git commit -m "docs(casascore-riesgos): plan y registro de riesgos [fecha prevista 22/09/2026]"
git push -u origin docs/casascore-riesgos
gh pr create --base main --reviewer xClaIDx --fill
\end{lstlisting}

\Needspace{11\baselineskip}
\noindent\textbf{PR 4 -- Plan de información y comunicaciones (22/09/2026)}
\begin{lstlisting}
git switch main && git pull
git switch -c docs/casascore-comunicaciones
cp -r ../paquete_viviana/docs/latex/casascore/12_plan_comunicaciones docs/latex/casascore/
git add docs/latex/casascore/12_plan_comunicaciones
git commit -m "docs(casascore-comunicaciones): plan de informacion y comunicaciones [fecha prevista 22/09/2026]"
git push -u origin docs/casascore-comunicaciones
gh pr create --base main --reviewer xClaIDx --fill
\end{lstlisting}

\Needspace{11\baselineskip}
\noindent\textbf{PR 5 -- Actas de los sprints 1 a 5 (16/10/2026)}
\begin{lstlisting}
git switch main && git pull
git switch -c docs/casascore-actas-sprint
cp -r ../paquete_viviana/docs/latex/casascore/15_actas_sprint docs/latex/casascore/
git add docs/latex/casascore/15_actas_sprint
git commit -m "docs(casascore-actas): actas de los sprints 1 a 5 [fecha prevista 16/10/2026]"
git push -u origin docs/casascore-actas-sprint
gh pr create --base main --reviewer xClaIDx --fill
\end{lstlisting}

\Needspace{11\baselineskip}
\noindent\textbf{PR 6 -- Plan de cierre y lecciones aprendidas (13/11/2026) -- súbelo el 13/11}
\begin{lstlisting}
git switch main && git pull
git switch -c docs/casascore-cierre
cp -r ../paquete_viviana/docs/latex/casascore/23_plan_cierre_lecciones docs/latex/casascore/
git add docs/latex/casascore/23_plan_cierre_lecciones
git commit -m "docs(casascore-cierre): plan de cierre y lecciones aprendidas [fecha prevista 13/11/2026]"
git push -u origin docs/casascore-cierre
gh pr create --base main --reviewer xClaIDx --fill
\end{lstlisting}

\noindent La opción \texttt{--fill} usa el mensaje del commit como título y descripción del Pull Request.

\Needspace{24\baselineskip}
\section{Manejo del Proyecto en el Navegador}
El tablero se usa desde \nolinkurl{https://github.com/users/xClaIDx/projects/1}. Escribe \nolinkurl{assignee:Viviana-11} en la barra de filtro para ver solo tus tarjetas.

\subsection{Qué haces con cada tarjeta}
{\small
\rowcolors{2}{white}{GrisFila}
\noindent\begin{tabularx}{\textwidth}{@{}P{4.4cm} C{1.9cm} L@{}}
\toprule
\rowcolor{AzulClaro}
\textbf{Tarjeta} & \textbf{Estado} & \textbf{Qué haces} \\
\midrule
Subir mis documentos a ProyectoTics (6 PR) & Por hacer & Arrástrala a \textit{En Progreso} al empezar. Ábrela y marca la casilla de cada PR cuando lo crees. Pásala a \textit{Hecho} cuando Clyde fusione el PR 6. \\
Riesgos v1.1: R02 activado & En Revision & Clyde revisa la actualización del registro de riesgos; atiende sus observaciones en la misma tarjeta. \\
Convocatoria al piloto & Backlog & Desde el 11/11: mensajes en grupos de WhatsApp de estudiantes y afiches; registra a los inscritos. \\
Plan de cierre y lecciones aprendidas & Backlog & Revisa el documento antes del 13/11 y súbelo con el PR 6 ese día. \\
Políticas y conformación del equipo; Planes de soporte; Pruebas del motor (7 casos) & Hecho & Cuando se fusione el PR de cada documento, escribe en la tarjeta un comentario con el enlace del Pull Request. \\
\bottomrule
\end{tabularx}
}


\Needspace{12\baselineskip}
\subsection{Acciones en la tarjeta}
\begin{itemize}[leftmargin=*]
    \item \textbf{Cambiar de estado:} arrastra la tarjeta a otra columna, o ábrela y cambia el campo \textit{Status}.
    \item \textbf{Marcar un avance:} abre la tarjeta y marca la casilla de la lista en su descripción.
    \item \textbf{Registrar horas:} en el panel derecho de la tarjeta, suma el tiempo en \textit{Horas reales}, el mismo día y en decimales (media hora $=$ 0.5).
    \item \textbf{Dejar evidencia:} escribe un comentario abajo de la tarjeta con el enlace del Pull Request.
    \item \textbf{Avisar un bloqueo:} escribe el motivo como comentario y avisa a Clyde por WhatsApp el mismo día.
\end{itemize}

\Needspace{16\baselineskip}
\section{Comandos de Consulta Rápida}
{\small
\rowcolors{2}{white}{GrisFila}
\noindent\begin{tabularx}{\textwidth}{@{}P{5.6cm} L@{}}
\toprule
\rowcolor{AzulClaro}
\textbf{Comando} & \textbf{Para qué sirve} \\
\midrule
\texttt{pwd} & Ver en qué carpeta estás; debe terminar en \texttt{ProyectoTics}. \\
\texttt{git status} & Ver en qué rama estás y qué archivos cambiaron. \\
\texttt{git branch} & Listar tus ramas; la actual tiene un asterisco. \\
\texttt{git log --oneline -5} & Ver los últimos 5 commits. \\
\texttt{git switch main \&\& git pull} & Volver a \texttt{main} y traer lo último de GitHub. \\
\texttt{gh pr status} & Ver el estado de tus Pull Requests. \\
\bottomrule
\end{tabularx}
}

\Needspace{18\baselineskip}
\section{Si Algo Sale Mal}
{\small
\rowcolors{2}{white}{GrisFila}
\noindent\begin{tabularx}{\textwidth}{@{}P{5.2cm} L@{}}
\toprule
\rowcolor{AzulClaro}
\textbf{Mensaje o situación} & \textbf{Solución} \\
\midrule
\texttt{fatal: not a git repository} & No estás dentro de la carpeta; escribe \nolinkurl{cd ~/Desktop/ProyectoTics}. \\
\texttt{Permission denied} o error 403 al hacer \texttt{push} & Acepta la invitación al repositorio y vuelve a ejecutar \texttt{gh auth login}. \\
\texttt{Author identity unknown} & Falta configurar tu nombre y correo (sección 3.2). \\
\texttt{cp: cannot stat} & La carpeta \texttt{paquete\_viviana} no está en el Escritorio, al lado de \texttt{ProyectoTics}. \\
\texttt{pathspec did not match} & La ruta está mal escrita; usa la tecla Tab para completarla. \\
\texttt{rejected (fetch first)} & Ejecuta \texttt{git pull} en tu rama y vuelve a hacer \texttt{git push}. \\
Hiciste un commit estando en \texttt{main} & No hagas \texttt{push}; avisa a Clyde y él te indica cómo moverlo a una rama. \\
\bottomrule
\end{tabularx}
}

\end{document}
